\section{Introduction}\label{sec:intro}

Video Large Language Models (VideoLLMs)~\citep{zhang2023video, maaz2024video, lin2024video, ren2024timechat, li2024llava} have recently emerged as a core paradigm for video understanding by integrating visual perception with language-based reasoning and response. These models have demonstrated strong performance across a wide range of tasks, including video captioning~\citep{yang2023vid2seq,wang2024tarsier,abdar2024review}, video question answering~\citep{li2024mvbench,wu2024longvideobench,fu2025video,song2025moviechat+}, and temporal grounding~\citep{lin2023univtg,qu2024chatvtg}. They also support diverse real-world applications, such as public safety monitoring~\citep{cournan2016improving,ahn2023safefac} and intelligent transportation~\citep{wan2020intelligent,vishal2024eyes,zhang2025language}. These advances highlight the potential of VideoLLMs to become a foundational technology for next-generation intelligent visual systems.

Despite this progress, most existing VideoLLMs are still designed for offline settings~\citep{shen2024longvu,zhang2024llava,wang2025internvl3}. In this paradigm, a complete video or a pre-collected segment is first buffered and then analyzed. Although this setup is well-suited for post hoc analysis, it limits the system's ability to respond promptly to ongoing events, making it less effective for applications such as real-time AI assistants~\citep{qian2025dispider,wang2025streambridge,seed2026vision}, live video narration~\citep{chen2025livecc,xu2025streamingvlm}, and interactive robotic systems~\citep{lynch2023interactive,fang2025robix}. To address this limitation, recent studies have begun to explore streaming VideoLLMs \citep{chen2024videollm,yao2025timechat,wang2025mmduet2,seed2026vision}. Systems built on streaming VideoLLMs can continuously process incoming video frames, incrementally update their understanding of the environment, and provide real-time responses. More importantly, they move beyond the traditional paradigm in which systems respond only after receiving explicit queries, enabling proactive responses based on information discovered from the video stream.

Recent research on streaming VideoLLMs has explored effective designs in architectural design, model training, and deployment efficiency~\citep{chen2024videollm,xu2025streamingvlm,qian2025dispider,wang2025mmduet2}. Despite these advances, existing approaches still face significant challenges. Current streaming VideoLLMs mainly fall into two categories: decoupled architectures and unified architectures. Decoupled architectures~\citep{qian2025dispider, wang2025streambridge} rely on two separately deployed models, where a trigger model determines whether the primary VideoLLM should respond. Because the trigger model does not share the same contextual state with the primary model and is typically much smaller in scale, the triggering accuracy and its consistency with response generation may be limited, leading to unstable system behavior. Unified architectures offer a higher performance upper bound, but these works~\citep{chen2024videollm,xu2025streamingvlm} are generally limited to captioning-style narration tasks and remain less effective for complex open-ended video question answering. Although recent full-duplex models~\citep{minicpmo45} integrate question answering into a single framework, they still lack sufficient robustness for long-duration streaming and may suffer from memory overflow or performance degradation during extended inference.

To address these issues, we introduce \textbf{AURA} (\textbf{A}lways-On \textbf{U}nderstanding and \textbf{R}eal-Time \textbf{A}ssistance), an end-to-end streaming visual interaction framework powered by a streaming VideoLLM. Through the co-design of data, algorithms, and systems, AURA achieves high stability, broad capabilities, and long-term endurance. Specifically, AURA tackles two fundamental challenges in streaming video understanding: 1) enabling a unified model to process video streams frame by frame and autonomously decide whether to remain silent or generate an appropriate response, and 2) stably handling unbounded video-text inputs over extended durations. To this end, we propose \textbf{Interactive Video Stream Context Management}, which integrates unbounded video frame streams and textual question-answer (QA) interactions into a limited context in an appropriate manner, enabling the model to perform silent observation and support diverse response modes. Based on the streaming response modes defined by this mechanism, we further develop a \textbf{Coarse-to-Fine Data Engine} with a five-stage pipeline to systematically construct training data for Real-Time QA, Proactive QA, and Multi-Response QA. In addition, to balance continuous silence with timely responses, we introduce a \textbf{Silent-Speech Balanced Loss} during training. Furthermore, we implement a \textbf{Real-time Streaming Inference Framework}, leveraging KV-cache reuse and related optimizations to enable efficient real-time inference over unbounded video streams.

We develop AURA on top of Qwen3-VL-8B-Instruct~\citep{bai2025qwen3} as a unified streaming visual assistant for real-time and real-world interaction over continuous video streams. After training, AURA achieves state-of-the-art results on streaming benchmarks while maintaining performance comparable to that of leading models on offline benchmarks. To further support practical real-time assistance, we integrate Automatic Speech Recognition (ASR) and Text-to-Speech (TTS) into the system and develop a functional demo application. With deployment optimizations, the system runs in real time at 2 FPS on two 80G accelerators. The primary contributions of our work are summarized as follows:

\begin{itemize}

\item We study streaming visual interaction with VideoLLMs as a setting that requires continuous observation, selective silence, and timely response over unbounded video streams. We analyze the limitations of existing offline, decoupled, and unified approaches in this setting.

\item We propose AURA, an end-to-end streaming visual interaction framework that enables a unified VideoLLM to process incoming frames continuously and support both real-time question answering and proactive responses.

\item We develop a co-designed pipeline spanning data construction, training, and deployment for stable long-horizon streaming, including Interactive Video Stream Context Management for context organization, a Coarse-to-Fine Data Engine for data construction, a Silent-Speech Balanced Loss for training, and a Real-time Streaming Inference Framework for deployment.

\item AURA achieves state-of-the-art results on streaming benchmarks while maintaining competitive performance on offline video understanding tasks. To demonstrate its practical usability, we build a real-time demo system with ASR and TTS, which runs at 2 FPS on two 80G accelerators.

\end{itemize}



